\section{Certificate escape analysis and stable observers}

\subsection{Three escape modes}

Every certificate-producing node receives one of three explicit modes.

\begin{description}[leftmargin=2.2cm,style=nextline]
\item[Local]
The public conclusion does not mention the certificate.  Use existential
elimination and erase it immediately.

\item[Observed]
The public conclusion depends only on \(O(c)\), and all accepted certificates
have the same observation, equivalent observations, or observations within a
specified diameter.  Expose the observation contract, not the raw code.

\item[Named]
The public theorem genuinely requires a particular numerical code.  Prove a
finite bound and existence, then define the least accepted code by bounded
search.
\end{description}

This classification is more than documentation.  It determines the required
proof obligations and prevents accidental global choice.

\subsection{Exact observers}

Let \(C(a,c)=1\) be a certificate checker and \(O(a,c)\) an observation term.
Assume pairwise stability:
\[
 C(a,c)=1\wedge C(a,c')=1
 \;\longrightarrow\;
 O(a,c)=O(a,c').
\]

\begin{theorem}[Exact observation erasure]
Suppose \(\exists c<B(a)\,C(a,c)=1\).  Then there exists an observation \(o\)
realized by an accepted certificate, and any two such observations are equal.
Any downstream formula congruent under equality of \(o\) can be proved without
naming a certificate.
\end{theorem}

\begin{proof}
Open the producer existential.  Introduce \(o=O(a,c)\) and prove the downstream
property from the local certificate.  Close the existential elimination; the
raw certificate disappears.  For uniqueness, open two certificates and apply
pairwise stability.  A downstream predicate is transported by equality using
the calculus's simultaneous Leibniz rule.  At no point is a selector over all
inputs constructed.
\end{proof}

If a concrete term for \(o\) must appear outside an existential, one may use the
least accepted certificate.  The point is that naming is an additional public
requirement, not a prerequisite for ordinary proof composition.

\subsection{Quotient observers without quotient types}

Exact equality is often too strong.  Bases differ by a change-of-basis matrix,
root lists differ by permutation, paths differ by stuttering, and local charts
differ on overlap.  Let \(E(x,y)=1\) be a reflected equivalence relation and
assume
\[
 C(a,c)=1\wedge C(a,c')=1
 \;\longrightarrow\;
 E(O(a,c),O(a,c'))=1.
\]
A consumer supplies an \(E\)-congruence theorem for its predicate.  This is
enough; no quotient type need be constructed inside arithmetic.

\begin{proposition}[Quotient observation principle]
If \(P(x)\) is stable under \(E\), then proving \(P(O(a,c))\) from one accepted
certificate proves an invariant public statement independent of the
certificate representation.
\end{proposition}

\begin{proof}
Given any other accepted \(c'\), observer stability gives \(E(O(c),O(c'))\).
The consumer's congruence implication transports \(P\) to \(O(c')\).  Because
the original certificate remains under existential elimination, the public
statement is the invariant property, not a chosen equivalence class object.
\end{proof}

This is particularly useful for unordered root clusters, finite sets represented
by duplicate-permitting lists, and geometric objects represented in overlapping
coordinate patches.

\subsection{Metric observers}

Approximation problems rarely have exact certificate independence.  Let
\(d(x,y)\) be a primitive-recursive rational pseudo-distance and let
\(\varepsilon(a,q)\) be an explicit error bound.  Assume
\[
 C(a,q,c)=1\wedge C(a,q,c')=1
 \;\longrightarrow\;
 d(O(a,q,c),O(a,q,c'))\le\varepsilon(a,q).
\]
Then any selected observation is representative up to the declared diameter.
If a consumer is \(L\)-Lipschitz, certificate choice contributes at most
\(L\varepsilon\) to the output error.

The key discipline is that \(\varepsilon\) is data, not a slogan such as
``sufficiently close.''  Its derivation may depend on separation, a modulus,
or a finite-cover certificate, each retained as a residual when unavailable.

\subsection{Independent certificates versus compatible paths}

There are two distinct ways to use a family of certificates.

\begin{enumerate}
\item If later stages require the actual state produced earlier, use relational
viability and a compatible replay.
\item If the final output is stable across all accepted certificates at each
query, select certificates independently.  Compatibility is unnecessary.
\end{enumerate}

The second method can be much cheaper.  For example, to approximate a symmetric
polynomial of all roots in a cluster, each precision may use a different list
ordering or a different isolating box.  Stable observation proves that the
symmetric output is consistent.  Constructing an infinite coherent branch
would solve a harder problem with no benefit to the requested theorem.

\subsection{Least naming}

Suppose a public API demands a code-valued function.  Given
\[
 \exists c<B(a)\;C(a,c)=1,
\]
define
\[
 \mu(a)=\min\{c<B(a):C(a,c)=1\}.
\]
The existing bounded-search theorem yields
\[
 \mu(a)<B(a),\qquad C(a,\mu(a))=1,
\]
and minimality.  This produces a canonical total primitive-recursive term
under the existence hypothesis.

A least code is representation dependent.  It should not be confused with a
mathematically canonical object.  Observer theorems ensure that public
invariant outputs do not depend on arbitrary details of the code ordering.

\subsection{Branch labels and discontinuity}

A uniform exact label for a root, eigenvector, or analytic branch may be
discontinuous at collisions or symmetries.  STRIVE does not force such a label.
It offers three safer interfaces:

\begin{itemize}
\item return the whole finite cluster;
\item return a symmetric or quotient-invariant observation;
\item return an approximate branch only under an explicit separation
certificate.
\end{itemize}

When separation is supplied, finite search can select the unique isolated
candidate and the observer becomes exact.  Without it, the residual accurately
records the missing mathematical resource.
